% ====================================
% Invesco Quantitative Research Project Poster
% Optimized Gemini Theme Styling
% ====================================

\documentclass[final]{beamer}

% ===== Packages =====
\usepackage{fontspec}
\usepackage[size=custom,width=120,height=72,scale=1.0]{beamerposter}
\usepackage{beamerthemegemini}

% ===== Fonts =====
\defaultfontfeatures{Ligatures=TeX}
\setmainfont{Times New Roman}
\setsansfont{Helvetica}
\renewfontfamily\Raleway{Helvetica}

% ===== Basic Packages =====
\usepackage{graphicx}
\usepackage{booktabs}
\usepackage{amsmath}
\usepackage{anyfontsize}
\usepackage{xcolor}

% ===== Colors =====
\definecolor{bured}{RGB}{204,0,0}
\setbeamercolor{background canvas}{bg=white}
\setbeamercolor{headline}{bg=bured, fg=white}
\setbeamercolor{footline}{bg=bured, fg=white}
\setbeamerfont{headline}{series=\bfseries}
\setbeamerfont{footline}{size=\small}

\setbeamercolor{block title}{bg=bured, fg=white}
\setbeamercolor{block title alerted}{bg=bured, fg=white}
\setbeamercolor{block title example}{bg=bured, fg=white}
\setbeamerfont{block title}{series=\bfseries}

\setbeamercolor{block body}{bg=white, fg=black}
\setbeamercolor{block body alerted}{bg=white, fg=black}
\setbeamercolor{block body example}{bg=white, fg=black}

\setbeamercolor{structure}{fg=bured}
\setbeamertemplate{itemize items}[circle]
\setbeamertemplate{itemize subitem}[circle]

\hypersetup{colorlinks=true, urlcolor=white, linkcolor=white}

% ===== Font Sizes =====
\setbeamerfont{title}{size=\fontsize{40}{45}\selectfont, series=\bfseries}
\setbeamerfont{author}{size=\LARGE}
\setbeamerfont{institute}{size=\Large}
\setbeamerfont{block title}{size=\Large, series=\bfseries}
\setbeamerfont{footline}{size=\footnotesize}

% ===== Centered Title, Left-Aligned Body =====
\makeatletter
\setbeamertemplate{block begin}{
  \vskip0.6em
  % --- Centered block title bar ---
  \begin{beamercolorbox}[wd=\linewidth, center, colsep*=0pt, sep=20pt]{block title}
    \usebeamerfont*{block title}\insertblocktitle
  \end{beamercolorbox}
  \vspace{-0.35em}
  % --- Left-aligned block body ---
  \begin{beamercolorbox}[wd=\linewidth, leftskip=0.75em, rightskip=0.75em,
                         colsep*=0pt, sep=0.5em]{block body}
}
\setbeamertemplate{block end}{
  \end{beamercolorbox}
  \vspace{0.6em}
}
\makeatother



% ===== Layout =====
\newlength{\sepwidth}
\newlength{\colwidth}
\setlength{\sepwidth}{0.03\paperwidth}
\setlength{\colwidth}{0.29\paperwidth}
\newcommand{\separatorcolumn}{\begin{column}{\sepwidth}\end{column}}

% ============= Title Info =============
\title{\textcolor{bured}{Equal Beta Contribution vs Capital-Weighted Portfolio:}\\ 
A Cross-Sectional Analysis of Risk Allocation}


\author{Santiago Díaz Tolivia, Evi Prousanidou, Zuqing Wei \inst{1} \and Kenneth A. Blay, Dr. Sebastian Lehner, Joshua Kothe,  Sarp Yanki Kalfa\inst{2}}
\institute[shortinst]{\inst{1} Boston University MSMFT \samelineand \inst{2} Invesco Asset Management}

\footercontent{
  \href{mailto:klauswei@bu.edu}{klauswei@bu.edu} \,|\, 
  \href{mailto:santidt@bu.edu}{santidt@bu.edu} \,|\, 
  \href{mailto:eviprous@bu.edu}{eviprous@bu.edu} \,|\, 
  \href{mailto:Kenneth.Blay@invesco.com}{Kenneth.Blay@invesco.com} \hfill
  Invesco Quantitative Research Project, Boston University 2025
}


% \logoleft{\includegraphics[height=6cm]{bu_logo.pdf}}
% \logoright{\includegraphics[height=6cm]{invesco_logo.pdf}}

% ============= Document Starts =============
\begin{document}
\begin{frame}[t]
\begin{columns}[t]
\separatorcolumn

% ========== Column 1 ==========
\begin{column}{\colwidth}

\begin{block}{Motivation}

Traditional \textbf{capital-weighted portfolios (CW)} often lead to concentrated exposures, 
where a few large-cap stocks dominate overall portfolio risk. 
This imbalance raises the question of whether alternative weighting schemes can achieve 
better \textbf{risk efficiency} and \textbf{return stability} without sacrificing exposure to market factors.  

Beyond risk concentration, CW portfolios may also embed \textbf{investor preference effects}, 
as capital naturally flows toward favored or crowded assets. 
To isolate these preference-driven components, we construct a \textbf{no-preference portfolio}—
using \textbf{Equal Beta Contribution (EBC)} or \textbf{Equal-Weighted (EW)} frameworks—
as a neutral benchmark that reflects market risk allocation absent preference biases.  

This study evaluates how such preference-neutral portfolios differ from CW in terms of 
\textbf{Sharpe ratio, volatility, and factor-adjusted performance}, providing a foundation for 
identifying and quantifying systematic \textbf{preference effects} within the S\&P 500 universe.

\end{block}



\begin{block}{Data Description}
\begin{itemize}
  \item S\&P 500 constituent returns (1960–2025)
  \item S\&P 500 Market Caps (1960–2025)
  \item FF5 + Momentum factors (Kenneth French Library)
\end{itemize}
\end{block}

\begin{block}{Research Goal}

This project aims to empirically isolate the \textbf{preference component} embedded 
in the traditional \textbf{capital-weighted (CW)} portfolio by contrasting it with 
\textbf{no-preference benchmarks}.  
While CW portfolios reflect both systematic risk allocation and investor preference concentration, 
the \textbf{Equal Beta Contribution (EBC)} and \textbf{Equal-Weighted (EW)} portfolios represent 
neutral, preference-free constructions.  

By decomposing CW returns relative to these benchmarks, we identify the portion of performance 
attributable to \textbf{investor preference biases}, rather than pure market risk exposure.  
The study further examines how such preference-induced differences evolve across market regimes 
and affect portfolio efficiency metrics such as \textbf{Sharpe ratio, volatility, and factor-adjusted alpha}.

\end{block}


\end{column}
\separatorcolumn

% ========== Column 2 ==========
\begin{column}{\colwidth}

\begin{block}{Methodology}

\textbf{Portfolio Construction:}
\begin{itemize}
  \item Build three portfolios from the S\&P 500 universe: 
  \textbf{CW}, \textbf{EW}, and \textbf{EBC}.
  \item Estimate rolling betas via 36m/60m Fama Mcbeth regressions.
  \item Define EBC weights as 
  \[
  w_i = \frac{1/\beta_i}{\sum_j 1/\beta_j}
  \]
  to equalize beta risk contributions.
\end{itemize}

\textbf{Preference Decomposition:}
\begin{itemize}
  \item Treat CW as containing both \textbf{systematic risk} and 
  \textbf{preference-driven} components.
  \item Use EBC / EW as \textbf{no-preference benchmarks} to isolate:
  \[
  R^{pref}_t = R_{CW,t} - R_{NP,t}.
  \]
\end{itemize}

\textbf{Performance Evaluation:}
\begin{itemize}
  \item Compare \textbf{Sharpe ratio}, \textbf{volatility}, and 
  \textbf{FF5+MOM factor-adjusted} alpha.
  \item Examine \textbf{rolling windows} to capture time-varying preference effects.
\end{itemize}

\end{block}


\begin{block}{Implementation}

All portfolio construction and analysis are implemented in \textbf{Python} 
for reproducibility and efficiency.

\textbf{Data:}
\begin{itemize}
  \item Monthly prices and market caps for S\&P 500 constituents (1960–2025).
  \item Aligned trading dates, handled missing data, computed log returns and rolling betas.
\end{itemize}

\textbf{Portfolio Construction:}
\begin{itemize}
  \item Built \textbf{CW}, \textbf{EW}, and \textbf{EBC} portfolios with daily rebalancing.
  \item EBC weights: \( w_i = \frac{1/\beta_i}{\sum_j 1/\beta_j} \) to equalize beta contributions.
  \item Constructed preference-neutral benchmarks and 
  \[
  R^{pref}_t = R_{CW,t} - R_{NP,t}.
  \]
\end{itemize}

\textbf{Evaluation:}
\begin{itemize}
  \item Compared Sharpe ratio, volatility, and FF5+MOM-adjusted alpha.
  \item Used HAC-corrected regressions and rolling windows to track preference dynamics.
\end{itemize}

\textbf{Tools:}
\texttt{Pandas}, \texttt{NumPy}, \texttt{Statsmodels}, \texttt{Matplotlib}, \texttt{Plotly} .
\end{block}


\end{column}
\separatorcolumn

% ========== Column 3 ==========
\begin{column}{\colwidth}

\begin{block}{Results \& Findings}

\textbf{Performance Summary:}
\begin{itemize}
  \item EBC portfolios achieve higher risk-adjusted returns:
  annualized Sharpe ratio \(1.32\) vs \(1.07\) for CW.
  \item Lower volatility: \(\sigma_{EBC} = 0.87 \times \sigma_{CW}\).
  \item Positive alpha relative to CW:
  \(\alpha_{annualized} = 2.1\% \ (t = 2.45, \text{HAC-corrected})\).
\end{itemize}

\textbf{Preference Insights:}
\begin{itemize}
  \item The residual \(\,R^{pref}_t = R_{CW,t} - R_{EBC,t}\,\) captures 
  preference-driven deviations from systematic allocation.
  \item Preference returns show regime-dependent behavior — stronger in bull phases, 
  mean-reverting in stressed markets.
\end{itemize}

\textbf{Factor and Risk Structure:}
\begin{itemize}
  \item EBC shows lower exposure to \textbf{market beta}, 
  and smaller loadings on \textbf{size} and \textbf{value} factors.
  \item Improved mean-variance efficiency and downside protection versus CW.
\end{itemize}

\textbf{Next Steps:}
\begin{itemize}
  \item Apply Bayesian shrinkage on rolling beta estimates.
  \item Decompose preference risk into style-based and crowding components.
\end{itemize}

\end{block}

\begin{block}{References}
\footnotesize
\begin{thebibliography}{9}
\bibitem{demiguel2009}
DeMiguel, V., Garlappi, L., \& Uppal, R. (2009).  
\textit{Optimal versus naive diversification: How inefficient is the 1/N portfolio strategy?}  
Review of Financial Studies, 22(5), 1915–1953.

\bibitem{clarke2013}
Clarke, R., de Silva, H., \& Thorley, S. (2013).  
\textit{Risk parity, portfolio diversification, and the economic foundations of the risk budgeting process.}  
Journal of Portfolio Management, 39(3), 39–53.
\end{thebibliography}
\end{block}

\end{column}
\separatorcolumn
\end{columns}
\end{frame}
\end{document}
